\documentclass{report}
\usepackage[utf8]{inputenc}
\usepackage[T1]{fontenc}
\usepackage{geometry}
\usepackage{graphicx}
\usepackage{float}
\usepackage{hyperref}
\usepackage[english,indonesian]{babel}

\title{Optimasi Pemrosesan Template LaTeX Menggunakan Pendekatan Berbasis AI}
\author{Academic Support Agent}
\date{\today}

\begin{document}
\maketitle

\begin{abstract}

\end{abstract}

\chapter{PENDAHULUAN}


\chapter{TINJAUAN PUSTAKA}


% Original content (commented out):
% \documentclass[12pt]{report}
\usepackage{blindtext}
\usepackage[T1]{fontenc}
\usepackage{fontspec}                            %Cambiar el tipo de letra
\usepackage[a4paper, margin=25.4mm]{geometry}    % Medidas de la hoja
\usepackage{graphicx}                            % Graficos
\usepackage[font=small,labelfont=bf]{caption}    % Poner captions en las imagenes
\usepackage{subcaption}                          % Poner subcaptions en las imagenes
\usepackage{hyperref}                            % Paquete para escribir links
\usepackage{multirow}
\usepackage{minted}                              % Paquete para escribir codigo de programacion
\usepackage{listings}                            % Paquete para escribir codigo de programacion
\usepackage{amsmath}                             % Paquete para matematicas
\usepackage{xcolor}                              % Access the named colour LightGray
\usepackage{indentfirst}                         % Identation o sangrado

\definecolor{LightGray}{gray}{0.9}
\usepackage[
backend=biber,
style=phys,                                     % Estilo de citado
sorting=ynt
]{biblatex}
\addbibresource{References.bib}

\setmainfont{Times New Roman}                   %Elige la fuente
\leftmark
\usepackage{fancyhdr}                           %import package fancy
\pagestyle{fancy}                               %use package fancy
\fancyhead{}                                    %clear head 
%\fancyfoot{}                                   %clear foot
\setlength{\headheight}{15.1pt}                   %size of the head and foot
\fancyhead[C]{\textbf{\leftmark}}               %head of the document
%\fancyfoot[R]{\thepage}                        %foot

\begin{document}
\begin{titlepage}
\centering
{\bfseries\fontsize{14pt}{28pt}\selectfont PONTIFICIA UNIVERSIDAD CATÓLICA DEL PERÚ \par}
{\bfseries\fontsize{14pt}{28pt}\selectfont FACULTAD DE CIENCIAS E INGENIERÍA \par}

\vfill

{\includegraphics[width=1\textwidth]{Images/logo_pucp.png}\par}

{\bfseries\fontsize{14pt}{28pt}\selectfont Title \par}

{\bfseries\fontsize{14pt}{28pt}\selectfont Trabajo de investigación para obtener el grado académico de Licenciado en Ciencias con mención en Física \par}

\vfill

{\bfseries\fontsize{14pt}{28pt}\selectfont AUTOR \par}
{\fontsize{14pt}{28pt}\selectfont You \par}

\vfill

{\bfseries\fontsize{14pt}{28pt}\selectfont ASESOR \par}
{\fontsize{14pt}{28pt}\selectfont Proffesor \par}

\vfill

{\fontsize{14pt}{28pt}\selectfont Lima, Agosto, 2023 \par}
\end{titlepage}
\pagenumbering{roman}
\setlength{\parindent}{1cm}            %Especifica que es 1cm de identado
\fontsize{12}{24}\selectfont           %cambia el tamano de letra y espaciado

\include{Chapters/0Abstract}
\include{Chapters/0Acknowledgements}
\tableofcontents
\clearpage
\pagenumbering{arabic}
\include{Chapters/1Introduction}
\include{Chapters/2.1Chapter1}
\include{Chapters/3Conclusions}

\renewcommand*{\bibfont}{\fontsize{10}{10}\selectfont} %change the size of the bibliography font
\printbibliography %Prints bibliography
\appendix
\include{Chapters/4Appendix}

\end{document}

\end{document}